\section{Nástroj BitcoinTxSubgraphBuilder}
\label{sec:safar-tool}

Šafář a spoluautori navrhli otvorený nástroj
\texttt{BitcoinTxSubgraphBuilder}, ktorý z verejných bitcoinových
transakcií zostaví vysvetliteľný dátový súbor. Pri každom príznaku
uchováva jeho význam, postup výpočtu a čas, ku ktorému sa hodnota
vzťahuje \cite{safar2026elliptic,bitcoinSubgraphBuilder}.

Rozbor vychádza zo zdrojového kódu revízie
\texttt{8e01503c0f9f} z 9. februára 2026; úplný identifikátor je uvedený
v bibliografickom zázname repozitára. Jednotlivé podkapitoly sledujú
vstupy, výpočty a uložené výstupy programov. Plný beh vyžaduje vlastný
Bitcoin Core, PostgreSQL, Blockbook, .NET a sekvenčné spracovanie
blockchainu. Praktická časť preto používa menší rozsah; architektúra
pôvodného nástroja určuje kontrolované riziká.

\subsection{Technické pojmy}

\begin{description}
  \item[API] Dohodnutý spôsob, ako si dva programy posielajú požiadavky
  a odpovede. Blockbook API napríklad vráti údaje transakcie podľa
  verejného \texttt{txid}.
  \item[RPC] Rozhranie na volanie funkcie iného programu. Nástroj cez
  Bitcoin Core RPC žiada úplný uzol o bloky a ich transakcie.
  \item[PostgreSQL] Databáza, v ktorej nástroj priebežne uchováva
  transakcie, triedy, vzdialenosti a odpovede Blockbook.
  \item[JSON] Textový formát odpovede API. Obsahuje pomenované polia,
  napríklad \texttt{vin}, \texttt{vout}, \texttt{fees} a
  \texttt{blockHeight}.
  \item[Preskok] Jedna hrana v grafe. Transakcia vo vzdialenosti dva
  preskoky je dostupná cez dve po sebe idúce hrany.
  \item[Parser blockchainu] Program, ktorý číta blok po bloku a mení
  surové transakcie na priebežné štatistiky.
  \item[Kontrolný bod] Súbor s rozpracovaným stavom. Po prerušení možno
  pokračovať od uloženej výšky bloku namiesto opakovania celého výpočtu.
  \item[Stav UTXO] Mapa ešte neminutých výstupov. Umožňuje pri spracovaní
  vstupu nájsť adresu a hodnotu staršieho výstupu.
  \item[Predspracovanie] Posledná úprava dát pred učením, napríklad
  normalizácia čísel a vytvorenie masiek pre tréning, validáciu a test.
\end{description}

\subsection{Čo nástroj rieši a čo potrebuje na vstupe}

Nástroj prijíma zoznam verejných bitcoinových \texttt{txid} s triedami.
Pri práci s Elliptic preto musí existovať samostatné mapovanie
\texttt{txId}$\rightarrow$\texttt{txid}. Anonymizovaný \texttt{txId}
ako vstup nestačí. Nástroj mapovanie nehľadá a neidentifikuje vlastníkov
adries. Spracovanie začína od verejných transakcií, ktoré poskytol
používateľ \cite{safar2026elliptic,bitcoinSubgraphBuilder}.

Obrázok~\ref{fig:safar-tool-architecture} rozdeľuje spracovanie na tri
vrstvy. Modrá časť zhromaždí transakcie a ich grafové okolie. Oranžová
časť prejde blockchain a vytvorí časovo správne štatistiky adries.
Zelená časť spojí oba zdroje do tabuliek použiteľných pri učení.

\begin{figure}[p]
  \centering
  \begin{tikzpicture}[
      >=Latex,
      every node/.style={font=\small,align=center},
      phase/.style={draw,rounded corners=2mm,minimum width=43mm,minimum height=12mm,line width=.7pt},
      store/.style={draw,cylinder,shape border rotate=90,aspect=.23,minimum width=34mm,minimum height=14mm,line width=.7pt},
      flowarrow/.style={-{Latex[length=2.4mm]},line width=.8pt},
      note/.style={font=\scriptsize,align=left,text width=42mm}
    ]
    \node[phase,fill=blue!9] (seed) at (0,0) {označené verejné\\\texttt{txid}};
    \node[phase,fill=blue!9] (import) at (0,-1.8) {Python: import\\počiatočných transakcií};
    \node[phase,fill=blue!9] (expand) at (0,-3.6) {Python: Blockbook\\rozšírenie na $d$ preskokov};
    \node[store,fill=blue!5] (postgres) at (0,-5.6) {PostgreSQL\\transakcie + JSON};

    \node[phase,fill=orange!12] (addresses) at (6.1,-1.8) {C\#: výber adries\\a cieľových výšok};
    \node[phase,fill=orange!12] (scan) at (6.1,-3.6) {C\#: Bitcoin Core\\blok po bloku};
    \node[store,fill=orange!7] (stats) at (6.1,-5.6) {štatistiky adries\\pri každej výške};

    \node[phase,fill=green!10] (construct) at (3.05,-7.7) {C\#: konštrukcia príznakov\\a hrán};
    \node[phase,fill=green!10] (preprocess) at (3.05,-9.5) {Python: predspracovanie\\iba podľa tréningu};
    \node[phase,fill=green!16] (output) at (3.05,-11.3) {CSV + DGL graf\\nový transparentný dátový súbor};

    \draw[flowarrow] (seed) -- (import);
    \draw[flowarrow] (import) -- (expand);
    \draw[flowarrow] (expand) -- (postgres);
    \draw[flowarrow] (postgres.east) to[out=0,in=-90] (addresses.south);
    \draw[flowarrow] (addresses) -- (scan);
    \draw[flowarrow] (scan) -- (stats);
    \draw[flowarrow] (postgres.south) to[out=-90,in=180] (construct.west);
    \draw[flowarrow] (stats.south) to[out=-90,in=0] (construct.east);
    \draw[flowarrow] (construct) -- (preprocess);
    \draw[flowarrow] (preprocess) -- (output);

    \node[note,anchor=west] at (8.6,-2.7)
      {Blockbook poskytne pohodlné\\vyhľadanie konkrétnej transakcie.};
    \node[note,anchor=west] at (8.6,-4.5)
      {Bitcoin Core poskytne úplnú\\chronologickú históriu blokov.};
  \end{tikzpicture}
  \caption{Architektúra nástroja BitcoinTxSubgraphBuilder. Dve vetvy sa
  spoja až pri výpočte príznakov: metadata transakcií pochádzajú z
  Blockbook, časové štatistiky adries zo sekvenčného čítania Bitcoin
  Core. Vlastné spracovanie podľa \cite{safar2026elliptic,bitcoinSubgraphBuilder}.}
  \label{fig:safar-tool-architecture}
\end{figure}

\subsection{Python: vytvorenie transakčného podgrafu}

Prvý program, \texttt{import\_elliptic\_to\_postgres.py}, načíta známe
verejné hashe, triedu a časový krok. Označené transakcie dostanú
vzdialenosť nula. Databázový riadok ďalej obsahuje pole
\texttt{distance}, pôvod v Elliptic a pole \texttt{data}, do ktorého sa
neskôr uloží odpoveď Blockbook.

Druhý program, \texttt{expand\_transaction\_network.py}, buduje okolie
po vrstvách. Z poľa \texttt{vin} číta staršie transakcie, ktorých
výstupy boli minuté. Z poľa \texttt{vout} číta \texttt{spentTxId}, teda
novšie transakcie, ktoré už daný výstup minuli. Okolie sa preto
rozširuje dozadu aj dopredu. Po každej vrstve program stiahne chýbajúce
JSON odpovede a cyklus opakuje, kým vo vzdialenosti najviac $d$ nič
nechýba.

Ak $S$ označuje množinu počiatočných označených transakcií a
$N^{-}(V)$ a $N^{+}(V)$ ich predchodcov a nasledovníkov, vrstvy možno
zapísať takto:
\begin{equation}
  V_0=S,\qquad
  V_k=V_{k-1}\cup N^{-}(V_{k-1})\cup N^{+}(V_{k-1}),
  \quad k=1,\ldots,d.
  \label{eq:hop-expansion}
\end{equation}
Vzdialenosť vrcholu $v$ je dĺžka najkratšej cesty od niektorého
počiatočného vrcholu:
\begin{equation}
  \operatorname{dist}(v,S)=\min_{s\in S}\operatorname{dist}(s,v).
\end{equation}
V praxi sa smer pri meraní okolia ignoruje, hoci vo výslednom grafe sa
smer toku zachováva.

\begin{figure}[htbp]
  \centering
  \begin{tikzpicture}[
      >=Latex,
      every node/.style={font=\small,align=center},
      tx/.style={draw,circle,minimum size=10mm,line width=.7pt},
      edge/.style={-{Latex[length=2mm]},line width=.75pt},
      ring/.style={draw,dashed,rounded corners=3mm,line width=.7pt}
    ]
    \node[ring,minimum width=28mm,minimum height=23mm,fill=green!5] (r0) at (0,0) {};
    \node[ring,minimum width=67mm,minimum height=44mm,fill=blue!3] (r1) at (0,0) {};
    \node[ring,minimum width=105mm,minimum height=64mm,fill=orange!3] (r2) at (0,0) {};
    \node[tx,fill=green!18] (seed) at (0,0) {$S$};
    \node[tx,fill=blue!15] (a) at (-2.1,.9) {$A$};
    \node[tx,fill=blue!15] (b) at (2.1,-.9) {$B$};
    \node[tx,fill=orange!18] (c) at (-4,.3) {$C$};
    \node[tx,fill=orange!18] (d) at (4,1.2) {$D$};
    \node[tx,fill=orange!18] (e) at (3.8,-1.8) {$E$};
    \draw[edge] (c) -- (a);
    \draw[edge] (a) -- (seed);
    \draw[edge] (seed) -- (b);
    \draw[edge] (b) -- (d);
    \draw[edge] (b) -- (e);
    \node[font=\scriptsize,anchor=west] at (5.4,1.4) {$d=2$};
    \node[font=\scriptsize,anchor=west] at (3.6,2.6) {$d=1$};
    \node[font=\scriptsize,anchor=west] at (1.1,1.0) {$d=0$};
  \end{tikzpicture}
  \caption{Rozšírenie podgrafu na dva preskoky. Vrcholy $A$ a $B$ sú
  vzdialené jeden preskok od označenej transakcie $S$; vrcholy $C$,
  $D$ a $E$ dva preskoky. Vlastné spracovanie.}
  \label{fig:hop-expansion}
\end{figure}

Výhoda pevnej vzdialenosti je priamo spojená s grafovým modelom. Sieť
s $L$ vrstvami prijíma správu najviac z $L$ preskokov. Ak dátový súbor
obsahuje okolie do vzdialenosti $d\geq L$, hranica vzorky neodreže
informáciu, ktorú by daný model vedel použiť. Voľba väčšieho $d$
zvyšuje počet vrcholov a pridáva neoznačené kontextové transakcie.

\subsection{C\#: časovo správna história adries}

Samotná odpoveď transakcie neobsahuje úplnú históriu každej adresy v
presnom minulom okamihu. Nástroj preto najprv príkazom
\texttt{extract-addresses} vyberie všetky adresy z podgrafu a ku každej
uloží výšku bloku, pri ktorej bude štatistika potrebná. Príkaz
\texttt{collect} potom číta bloky cez Bitcoin Core RPC od začiatku
reťazca po zadanú cieľovú výšku.

Pri výstupe transakcie analyzátor pridá nový záznam do vlastnej mapy UTXO.
Pri neskoršom vstupe ten istý záznam nájde, prečíta jeho adresu a
hodnotu a z mapy ho odstráni. Tým určí pôvod hodnoty vstupu priamo z
predchádzajúceho výstupu. Počítadlá sa aktualizujú pre adresy potrebné
vo vybranom podgrafe; analyzátor pritom prechádza úplnú chronológiu
blokov.

Pre adresu $a$ a výšku bloku $h$ nástroj uchováva najmä tieto hodnoty:
\begin{align}
  R_a(h) &= \sum_{o:\,\operatorname{addr}(o)=a,\,h(o)\leq h}\operatorname{value}(o),
  &&\text{celková prijatá hodnota},\\
  S_a(h) &= \sum_{i:\,\operatorname{addr}(i)=a,\,h(i)\leq h}\operatorname{value}(i),
  &&\text{celková odoslaná hodnota},\\
  B_a(h) &= R_a(h)-S_a(h),
  &&\text{zostatok},\\
  L_a(h) &= t_a^{\mathrm{last}}(h)-t_a^{\mathrm{first}}(h),
  &&\text{životnosť adresy}.
  \label{eq:address-stats}
\end{align}
Ďalej sa počíta počet transakcií, počet vytvorených výstupov, počet
minutých výstupov a počty prichádzajúcich a odchádzajúcich transakcií.
Hodnota pri $h$ obsahuje iba udalosti s výškou najviac $h$. Transakcia
z budúceho bloku $h+1000$ teda nemôže zmeniť príznak staršej
transakcie pri $h$.

\begin{figure}[htbp]
  \centering
  \begin{tikzpicture}[
      >=Latex,
      every node/.style={font=\small,align=center},
      event/.style={draw,circle,minimum size=9mm,line width=.7pt},
      snapshot/.style={draw,rounded corners=2mm,fill=green!10,minimum width=33mm,minimum height=11mm},
      timearrow/.style={-{Latex[length=2.2mm]},line width=.8pt}
    ]
    \draw[timearrow] (-.5,0) -- (11.2,0) node[right]{výška bloku};
    \node[event,fill=blue!10] (e1) at (1,0) {$h_1$};
    \node[event,fill=blue!10] (e2) at (4,0) {$h_2$};
    \node[event,fill=orange!15] (target) at (7,0) {$h$};
    \node[event,fill=gray!12] (future) at (10,0) {$h+q$};
    \node[snapshot] (state) at (7,-1.7) {príznak pri $h$\\používa $h_1,h_2,h$};
    \draw[timearrow] (e1.south) to[bend right=18] (state.west);
    \draw[timearrow] (e2.south) to[bend right=9] (state.north west);
    \draw[timearrow] (target.south) -- (state.north);
    \draw[dashed,line width=.8pt,red!70] (future.south) to[bend left=20]
      node[below,font=\scriptsize,text=red!70]{nepoužiť} (state.east);
  \end{tikzpicture}
  \caption{Časovo ohraničená štatistika adresy. Budúca udalosť sa do
  príznaku staršej transakcie nezapočíta. Vlastné spracovanie podľa
  \cite{safar2026elliptic,bitcoinSubgraphBuilder}.}
  \label{fig:time-scoped-address-stat}
\end{figure}

Spracovanie celého blockchainu trvá dlho. Implementácia preto sťahuje
bloky do obmedzeného kanála, spracúva ich v poradí a každých 50~000
blokov ukladá stav UTXO aj adresné štatistiky. Prerušený beh možno
obnoviť voľbou \texttt{--checkpoint}. Samostatný príkaz
\texttt{deduplicate} zlúči rovnaké štatistické snímky platné pre viac
cieľových výšok a zmenší výsledný súbor.

\subsection{Konštrukcia príznakov a hrán}

Príkaz \texttt{create-dataset} spojí tri zdroje: JSON transakcie z
PostgreSQL, časové štatistiky adries a hodinovú históriu ceny BTC.
Vytvorí súbor \texttt{tx\_features.csv} a súbor \texttt{edges.csv}.
Hrana sa uloží iba vtedy, keď sa v podgrafe nachádza zdrojová aj cieľová
transakcia. Nesie aj hodnotu minutého výstupu a výšky oboch blokov.

Tabuľka~\ref{tab:open-elliptic-features} uvádza všetky skupiny hodnôt,
ktoré konštruktor počíta. Štatistiky sa počítajú samostatne pre vstupy
a výstupy. Pri každej adresnej vlastnosti sa ukladajú minimum, maximum,
priemer, medián a rozpätie. \emph{Medián} je prostredná hodnota po
zoradení a \emph{rozpätie} je maximum mínus minimum.

\begin{table}[p]
  \centering
  \small
  \caption{Skupiny príznakov vytvárané nástrojom
  BitcoinTxSubgraphBuilder \cite{bitcoinSubgraphBuilder}.}
  \label{tab:open-elliptic-features}
  \begin{tabularx}{\textwidth}{lX}
    \toprule
    Skupina & Význam \\
    \midrule
    základ transakcie & výška bloku, hodnota a hodnota v USD, poplatok
      a poplatok v USD, veľkosť v bajtoch a virtuálnych bajtoch \\
    pomery & poplatok ku vstupnej hodnote, počet vstupov k počtu
      výstupov a počet vstupných transakcií k počtu výstupov \\
    typ transakcie & príznak coinbase a počet vstupných transakcií \\
    opakovanie adries & počet adries spoločných pre vstupy aj výstupy,
      podiel nových výstupných adries a podiel hodnoty smerujúcej na ne \\
    hodnoty vstupov/výstupov & počet, minimum, maximum, smerodajná
      odchýlka, priemer, medián a tie isté sumy po prepočte na USD \\
    podiely hodnôt & minimum, maximum, odchýlka, rozptyl a medián podielu
      jednej položky na celkovej vstupnej alebo výstupnej hodnote \\
    adresy & celkový a jedinečný počet adries a ich vzájomný pomer \\
    história adries & minimum, maximum, priemer, medián a rozpätie pre
      počet transakcií, prichádzajúce a odchádzajúce transakcie,
      zostatok, prijatú hodnotu, odoslanú hodnotu a životnosť \\
    \bottomrule
  \end{tabularx}
\end{table}

Tri odvodené veličiny možno zapísať priamo:
\begin{align}
  r_{\mathrm{fee}} &= \frac{\mathrm{fee}}{\mathrm{valueIn}},\\
  r_{\mathrm{addr}} &=
  \frac{|A_{\mathrm{out}}\setminus A_{\mathrm{in}}|}{|A_{\mathrm{out}}|},\\
  r_{\mathrm{foreign\ value}} &=
  \frac{\sum_{o:\,A(o)\cap A_{\mathrm{in}}=\emptyset}\mathrm{value}(o)}
  {\sum_o\mathrm{value}(o)}.
  \label{eq:open-elliptic-ratios}
\end{align}
Prvý pomer ukazuje, aká časť vstupnej hodnoty tvorí poplatok. Druhý
ukazuje podiel výstupných adries, ktoré sa nenachádzajú medzi vstupmi.
Tretí ukazuje podiel hodnoty odoslanej na takéto adresy. Určenie
vráteného zvyšku vyžaduje ďalšie údaje alebo heuristiku.

Pri kontrole revízie uvedenej na začiatku kapitoly sa našiel aj detail,
ktorý treba pri reprodukcii otestovať: výraz pre
\texttt{VinVoutRatio} delí dve celé čísla pred uložením do typu
\texttt{double}. V jazyku C\# sa tak výsledok najprv skráti na celé
číslo. Ostatný podobný pomer \texttt{TxInVoutRatio} obsahuje explicitný
prevod na \texttt{double}. Táto práca preto nepreberá hodnotu
\texttt{VinVoutRatio} bez samostatného testu a uvádza použitú revíziu
kódu, aby bol výsledok dohľadateľný.

\subsection{Predspracovanie pre strojové učenie}

Surový výstup sa predvolene nenormalizuje. Výskumník tak vidí pôvodné
satoshi, počty a časy a môže ich nezávisle skontrolovať. Voliteľný
skript \texttt{open\_\allowbreak{}elliptic\_\allowbreak{}preprocessing.py}
vytvorí formát podobný Elliptic. Kroky 1--29 používa na tréning, 30--34
na validáciu a 35--49 na test.

Pre každý numerický príznak $j$ sa priemer a smerodajná odchýlka
vypočítajú iba z transakcií dostupných na konci tréningovej časti:
\begin{equation}
  z_{ij}=\frac{x_{ij}-\mu_j^{\mathrm{train}}}
  {\sigma_j^{\mathrm{train}}}.
  \label{eq:safar-train-scaling}
\end{equation}
Rovnaké $\mu_j^{\mathrm{train}}$ a $\sigma_j^{\mathrm{train}}$ sa potom
použijú na validáciu aj test. Skript uloží aj objekt
\texttt{scaler.pkl} s parametrami škálovania, aby sa rovnaká transformácia dala použiť na novej
transakcii. Pri hranách pridá opačný smer pre učenie GNN. Hrany, ktoré
by po obrátení prenášali správu z neskoršej časti do skoršej, odstráni.

\subsection{Presný postup spustenia}

Tabuľka~\ref{tab:safar-commands} oddeľuje príkazy podľa vstupu a
výstupu. Každý krok zodpovedá jednej overiteľnej transformácii dát.

\begin{table}[p]
  \centering
  \scriptsize
  \caption{Postup nástroja podľa dokumentácie a zdrojového kódu
  \cite{bitcoinSubgraphBuilder}.}
  \label{tab:safar-commands}
  \begin{tabularx}{\textwidth}{lXX}
    \toprule
    Krok & Vstup a operácia & Výstup \\
    \midrule
    SQL tabuľky & vytvorenie tabuliek transakcií a adries v PostgreSQL
      & prázdna databázová schéma \\
    import & verejné \texttt{txid}, triedy a časové kroky
      & počiatočné vrcholy so vzdialenosťou 0 \\
    fetch/expand & Blockbook a maximálna vzdialenosť $d$
      & JSON všetkých transakcií podgrafu \\
    extract-addresses & vstupy a výstupy transakcií podgrafu
      & dvojice adresa--cieľová výška \\
    collect & bloky z Bitcoin Core od výšky 0 po cieľ
      & časové snímky štatistík a stav UTXO \\
    deduplicate & opakované rovnaké snímky jednej adresy
      & komprimované štatistiky pre viac výšok \\
    create-dataset & transakcie, štatistiky adries a história ceny
      & \texttt{tx\_features.csv} a \texttt{edges.csv} \\
    predspracovanie & surové CSV a hranice častí
      & normalizované CSV, triedy, hrany a parametre škálovania \\
    prevod do DGL & CSV vo formáte Elliptic
      & binárny graf s príznakmi, triedami a maskami \\
    \bottomrule
  \end{tabularx}
\end{table}

Minimálny logický sled začne importom a rozšírením grafu. Potom sa
extrahujú adresy, zozbierajú a skomprimujú ich štatistiky a napokon sa
vytvorí dátový súbor. Ak sa zmení
maximálna vzdialenosť $d$, zoznam adries alebo cieľová výška, nadväzujúce
výstupy treba vytvoriť znovu alebo jednoznačne oddeliť podľa konfigurácie.

\subsection{Opravené chyby a zostávajúce obmedzenia}

Nástroj zachováva pôvodný význam a jednotku príznakov, adresnú históriu
zastaví pri čase analyzovanej transakcie, normalizáciu učí na
tréningovej časti a grafové okolie vymedzuje
počtom preskokov. Takýto výstup možno vypočítať aj pre novú verejnú
transakciu, čo pri anonymizovaných stĺpcoch pôvodného Elliptic nebolo
možné \cite{safar2026elliptic}.

Nástroj nezaručí nezávislé označenia. Ak trénovacia a
testovacia množina zdieľajú adresu alebo tú istú entitu pod viacerými
adresami, model môže stále rozpoznávať už videný subjekt. Autori preto
odporúčajú vytvoriť nezávislý test, ktorý nezdieľa adresy ani adresy v
$d$-preskokových okoliach. Ani tento technický filter nedokazuje, že dve
odlišné adresy patria rôznym ľuďom.

Ďalším obmedzením je prevádzková náročnosť. Blockbook uľahčuje náhodný
prístup k transakciám; časovo presná história veľkého množstva adries
si vyžaduje sekvenčný prechod blockchainu. Výsledok tiež nemožno
publikovať ako opravený Elliptic, pretože licencia
CC BY-NC-ND 4.0 zakazuje šírenie odvodených upravených verzií. Preto je
potrebné publikovať zdrojový kód, konfiguráciu, verziu vstupov a súhrnné
výsledky. Chránený odvodený dátový súbor sa nezverejňuje
\cite{safar2026elliptic}.

Praktická časť tejto práce používa menší postup. Cez verejný Blockbook
získava šesť časových krokov, odpovede ukladá do SQLite a počíta 18
priamo overiteľných lokálnych príznakov. Bez sekvenčného analyzátora
Bitcoin Core nepočíta úplné časové štatistiky adries. Zvolený rozsah
oddeľuje reprodukovateľný bakalársky experiment od infraštruktúry
potrebnej na tvorbu celého nového dátového súboru.
